\section{Implications and Limits}\label{sec:agenda}
\subsection{A deployment evaluation checklist}
Before claiming that a defense protects a deployment, report: (1) the hidden target and protection relation; (2) the final receiver's actual fields and linkage permissions; (3) where protection ends and whether cover reaches action dispatch; (4) protected-data access, fitting procedure and a matched defended-test contrast where adaptation is claimed; (5) the workload and the generator shifts actually tested; and (6) inference, objective/policy utility and applicable cost bound to the same condition. Report missing entries as scoped evidence states. The checklist asks for evidence needed by the specific claim, not every measurement in every paper.

The three cases make this actionable. A content-aware downstream renderer calls for paired final-hop captures; surviving actions call for residual-feature prediction; executed dummy actions call for objective state/goal tests. Prefix-cache isolation needs workload reuse costs, while a semantic matcher needs a separate mechanism check. Task, topology, browser configuration and hardware shifts need their respective bridge tests. These follow from the source comparisons rather than from a global defense ranking.

\subsection{Research questions implied by the cases}
Five concrete questions follow from the decision procedure in Figure~\ref{fig:defense-decision}. (1) Across API, frontend and final receiver, which transformations preserve a protection bound, and which introduce additional target-dependent state? (2) Can a cover policy be certified goal- and policy-neutral before tool dispatch, while accounting for resource and timeout effects? (3) How much protected training data can each real observer obtain, and what is the resulting adaptation frontier under a fixed budget? (4) For isolation or equalization, how do workload reuse and the actual matching primitive change total service cost? (5) Which receiver-visible invariants survive changes in task, topology and backend without relying on private collector fields? These questions prioritize missing bridges over another heterogeneous attack ranking.

\subsection{Threats to the synthesis}
\paragraph{Selection and source binding}
The purposive collection supports source-specific deployment judgments rather than literature-wide prevalence. The completed finite update adds positive counterconditions and closer framework prior work; its 16 missing-abstract records and one inaccessible primary source remain material requests (S10). Historical acquisition gaps and four unbound defense quantities are separately retained (S1/S5). Case-level dependence is explicit in Table~\ref{tab:synthesis-decisions}; source deletion withdraws a source instance, while the new controlled tests retain their own independent local measurements.

\paragraph{Interpretive reliability}
Eligibility, result splitting and transfer judgments use a single AI-assisted coding workflow. Independent human agreement is unmeasured. The supplied locators, pairwise reasons and alternative dispositions make the interpretations auditable. The retained negative checklist comparison is the reason we position SCOPE as fields, not as a more effective review method.

\paragraph{Empirical reach}
The controlled checks bind receiver sizes and task states to actual interventions in a finite local system. They isolate placement through matched output/plan replay; they do not measure remote generation timing or a deployed adaptive agent population. The supplementary archive case addresses prediction under archive-receiver ordering, where matched-task evidence remains inconclusive. Original failed pilots and retail results are retained alongside the new controlled results.

\paragraph{Disclosure and materials}
The companion Supplementary Evidence contains source-linked tables, saved predictions, protocols and a local verifier, excluding private logs, conversations and redistributed third-party datasets. It is prepared for submission as accompanying material; the README provides the entry points. The authors remain responsible for AI-assisted judgments and submission disclosures.
